\begin{recipe}
[% 
    preparationtime = {\unit[30]{min}},
    bakingtime = {\unit[30]{min}},
    bakingtemperature = {\protect\bakingtemperature{
        topbottomheat=\unit[180]{\textcelcius}
    }},
    portion = {\portion{9}},
    source = {\protect\recipesource{https://lovingitvegan.com/vegan\-sticky\-toffee-pudding/\#recipe}{Loving It Vegan: Vegan Sticky Toffee Pudding}},
    calory = {444kcal | 4g Protein | 56g KH | 23g Fett},
    price = {1,25€},
    created = {07.06.2026},
    %rating = {1},
]
{Veganer Sticky Toffee Pudding}

\graph{% pictures
    big=pic/dessert/12_veganer-sticky-toffee-pudding
}

\introduction{%
Ein veganer Sticky Toffee Pudding mit Datteln, warmen Gewürzen und einer richtig klebrigen Toffee-Sauce. Absolut kein leichtes Dessert, aber das war hier auch wirklich nie der Plan.
}

\ingredients{
    \textbf{Pudding} & \\
    175 g & Entsteinte Medjool-Datteln\index{Obst!Datteln}\\
    360 ml & Sojamilch\index{Milchprodukte \& Alternativen!Sojamilch}\\
    1 TL & Vanilleextrakt\index{Backzutaten \& Süßes!Vanille}\\
    112 g & Vegane Butter\index{Milchprodukte \& Alternativen!Vegane Butter}\\
    100 g & Brauner Zucker\index{Backzutaten \& Süßes!Zucker}\\
    188 g & Weizenmehl\index{Getreide, Pasta \& Brot!Weizenmehl}\\
    3 TL & Backpulver\index{Backzutaten \& Süßes!Backpulver}\\
    1 TL & Salz\index{Kräuter \& Gewürze!Salz}\\
    \nicefrac{1}{2} TL & Muskat\index{Kräuter \& Gewürze!Muskat}\\
    1 TL & Gemahlener Ingwer\index{Gemüse \& Pilze!Ingwer}\\
    1 TL & Zimt\index{Kräuter \& Gewürze!Zimt}\\
    \textbf{Toffee-Sauce} & \\
    1 EL & Ahornsirup\index{Öle, Saucen \& Würzmittel!Ahornsirup}\\
    100 g & Dunkler brauner Zucker\index{Backzutaten \& Süßes!Zucker}\\
    112 g & Vegane Butter\index{Milchprodukte \& Alternativen!Vegane Butter}\\
    120 ml & Kokoscreme\index{Milchprodukte \& Alternativen!Kokoscreme}
}

\preparation{%
    \step%
    Backofen auf 180 °C Ober-/Unterhitze vorheizen und eine quadratische Backform einfetten.
    
    \step%
    Datteln grob hacken und zusammen mit Sojamilch und Vanilleextrakt in einen Topf geben. Bei mittlerer Hitze etwa 10--12 Minuten erhitzen und dabei regelmäßig rühren, bis die Datteln weich sind.
    
    \step%
    Vegane Butter und braunen Zucker zur warmen Dattelmasse geben und einrühren, bis beides geschmolzen ist. Vom Herd nehmen.
    
    \step%
    Mehl, Backpulver, Salz, Muskat, Ingwer und Zimt in einer Schüssel vermengen.
    
    \step%
    Die warme Dattelmasse zu den trockenen Zutaten geben und zu einem dicken Teig verrühren.
    
    \step%
    Teig in die vorbereitete Form geben, glatt streichen und etwa 30 Minuten backen, bis bei der Stäbchenprobe kein Teig mehr kleben bleibt.
    
    \step%
    Den Pudding etwa 10 Minuten abkühlen lassen.
    
    \step%
    Währenddessen für die Sauce Ahornsirup, braunen Zucker und vegane Butter in einem Topf erhitzen, bis die Mischung blubbert.
    
    \step%
    Vom Herd nehmen, Kokoscreme einrühren und zu einer glatten Sauce verrühren.
    
    \step%
    Pudding in Stücke schneiden und großzügig mit der warmen Toffee-Sauce servieren.
}

\suggestion[Servieren]{%
    Sehr gut mit veganem Vanilleeis, veganer Schlagsahne oder Custard. Die Nährwerte oben beziehen sich nur auf Pudding und Sauce, nicht auf zusätzliches Eis oder Sahne.
}

\hint{
    Die Sauce ist bewusst großzügig gerechnet. Also nicht geizig sein, sonst ist es am Ende nur Kuchen mit Identitätskrise.
}

\end{recipe}